\section*{Chapter \ref{ch:s1:merger-arbitrage} --- Merger Arbitrage}

\begin{solution}{exo:s1:merger-arbitrage:1}
Spread $40/38.60 - 1 = 3.63\%$; implied probability $(38.60 - 30)/(40 - 30) = 0.86$.
\end{solution}

\begin{solution}{exo:s1:merger-arbitrage:2}
$(1.0363)^{252/84} - 1 = 11.3\%$ a year.
\end{solution}

\begin{solution}{exo:s1:merger-arbitrage:3}
At \$42, inside the band, the consideration is $0.5 \times 42 = \$21$ and the hedge is 0.5 acquirer shares short. At \$36, below the band, the consideration is
fixed at $0.5 \times 40 = \$20$ and the hedge is zero: the value no longer depends on the acquirer.
\end{solution}

\begin{solution}{exo:s1:merger-arbitrage:4}
Expected profit $0.86 \times (40 - 38.60) - 0.14 \times (38.60 - 30) = 1.204 - 1.204 = 0$. To trade, the arbitrageur must believe the break probability is below
14\% (or the \omterm{def:s1:merger-arbitrage:spread}{break price} above \$30), after discounting and costs.
\end{solution}

\begin{solution}{exo:s1:merger-arbitrage:5}
Each deal's \omterm{def:s1:merger-arbitrage:spread}{break price} moves with the market, so a market fall marks every open position down at once, whether or not any deal breaks; and more deals break
when the market falls. Both are common to all deals and do not diversify.
\end{solution}

\begin{solution}{exo:s1:merger-arbitrage:6}
It extends the review by months and signals antitrust concern: the break probability rises and the time to completion lengthens, so the spread widens and its
annualised return falls for the same absolute spread.
\end{solution}

\begin{solution}{exo:s1:merger-arbitrage:7}
With breaks independent of the market the portfolio earns 4.30\% over the rate (Sharpe ratio 1.50): the clustering of breaks in downturns costs 1.5 points a
year. With spreads priced at the flat-market break rate it loses 0.47\% (Sharpe ratio $-0.17$): the premium in the base-case spread is what pays for the
crash-dependent breaks.
\end{solution}

\begin{solution}{exo:s1:merger-arbitrage:8}
Its beta near zero is an average: close to zero in rising and flat markets and higher in falling ones (0.20 against 0.09 here, larger in the historical
data), like a short put. It is short crash risk, and a Sharpe ratio built on months without crashes overstates its quality.
\end{solution}

\begin{solution}{pb:s1:merger-arbitrage:1}
\begin{enumerate}
\item Buying targets below the consideration (hedging stock deals) to earn the spread on completion; an offer made directly to shareholders; offer over price
minus one; the price if the deal fails.
\item 0.86.
\item Deals take months; the same spread over different horizons is a different return.
\item About fifty cash deals a year, premiums of 20--40\%, a median of 110 days, \omterm{def:s1:merger-arbitrage:spread}{break prices} moving with the market, spreads pricing a 12\% break
probability, true break probability 6\% plus 1.5 times the market's fall.
\item 6.76\% a year, 2.76\% over the rate, volatility 3.26\%, Sharpe ratio 0.85.
\item 88\% implied, 92.5\% realised.
\item 13.7\% when the market fell over the deal's life, 4.7\% when it rose.
\item 0.20 in months the market fell over 4\%, 0.09 otherwise.
\item Independent breaks earn 4.30\%: clustering costs 1.5 points; spreads priced at 6\% lose 0.47\%: the premium pays for the crash risk.
\item $-3.06\%$ in a month the market rose 2.8\%: a few deal-specific breaks together.
\item Short the exchange ratio of acquirer shares; with a collar, short its delta, which jumps at the band's edges.
\item Antitrust review, financing, shareholder votes, competing bids, material adverse changes.
\item \textbf{Named result.} The portfolio's beta is 0.20 in months when the market falls more than 4\% and 0.09 in other months; the spreads imply an 88\%
completion probability and 92.5\% of deals complete.
\item Returns correlated with the market in severe falls and not otherwise, like a short index put; about 4\% a year in excess returns after costs, 1963--1998.
\item A short out-of-the-money index put.
\item By the loss if it breaks, relative to the book.
\item At least 30 days (15 for cash \omterm{def:s1:merger-arbitrage:arb}{tender offers}) after complete filings, longer with a Second Request.
\item Deal-specific breaks; not market-dependent breaks and \omterm{def:s1:merger-arbitrage:spread}{break prices}.
\item After a fall, when spreads are wide, if the capital can bear further falls.
\item Insurance against deals breaking, which pays most in bad markets.
\end{enumerate}
\end{solution}

\begin{solution}{iq:s1:merger-arbitrage:1}
Buying the target of an announced acquisition below the offer (and shorting the acquirer in stock deals) to earn the spread when the deal closes. The return
is a premium for bearing break risk, much of which comes in market downturns.
\end{solution}

\begin{solution}{iq:s1:merger-arbitrage:2}
Short the collar's delta in acquirer shares: the exchange ratio inside the band, zero outside it where the value is fixed; rebalance as the acquirer's price
approaches the band's edges, where the delta jumps, or buy options that offset the collar's gamma.
\end{solution}

\begin{solution}{iq:s1:merger-arbitrage:3}
Its payoff is near zero exposure in normal markets and losses in crashes, when deals break and \omterm{def:s1:merger-arbitrage:spread}{break prices} fall together: small steady premiums with rare
large correlated losses, the profile of a short put.
\end{solution}

\begin{solution}{iq:s1:merger-arbitrage:4}
\omterm{def:s1:merger-arbitrage:spread}{Break price} from the undisturbed price moved by the market and the sector since the announcement, or from peers; break probability from deal features
(regulatory overlap, financing, hostility, vote), base rates of similar deals, and the market-implied probability as a prior.
\end{solution}

\begin{solution}{iq:s1:merger-arbitrage:5}
Rumours of regulatory trouble or financing problems, a market fall moving \omterm{def:s1:merger-arbitrage:spread}{break prices}, forced selling by other arbitrageurs, or a borrow problem in the
acquirer for a stock deal.
\end{solution}

\begin{solution}{iq:s1:merger-arbitrage:6}
With price $A = e^{-r\tau}(pK + (1 - p)B)$, the implied probability is $p = (Ae^{r\tau} - B)/(K - B)$. If the true probability is $\pi$, the expected payoff is
$\pi K + (1 - \pi)B$ and the expected excess return is $(\pi - p)(K - B)/(Ae^{r\tau})$: proportional to the gap between the true and implied probabilities
and to the spread between the offer and the \omterm{def:s1:merger-arbitrage:spread}{break price}.
\end{solution}
